\section{Settings, help and keyboard shortcuts}
\label{sec:settings}

The main window has three menus: \ui{File}, \ui{Settings} and \ui{Help}. This section covers the
\ui{Settings} and \ui{Help} menus, lists the \ui{File} menu with the sections that describe its
items, and gives the keyboard shortcuts. RAFTcorr remembers the theme and the language for the next start, in
the settings of your user account (Section~\ref{sec:overview-files}).

\subsection{The Settings menu}
\label{sec:settings-menu}

\begin{description}[leftmargin=0.6cm,style=nextline]
  \item[\menuthree{Settings}{Theme}{Dark} or \ui{Light}]
    The colour theme of the interface. \ui{Light} is the default. The change is immediate: every
    open window, the strain window included, is drawn again in the new colours, and the log keeps
    its lines.
  \item[\menu{Settings}{Language}]
    The language of the interface: English, Simplified Chinese, Traditional Chinese, Japanese,
    Korean, German, French or Spanish, each listed under its own name. RAFTcorr saves your choice
    and shows \ui{Language changed}, which asks you to restart: the windows already open keep
    their labels until then. The first time RAFTcorr starts, it takes the language of the
    operating system when it has a translation for it (Portuguese falls back to Spanish), and
    English otherwise.
\end{description}

Messages that come from the tracking engine itself, such as the tiling line and the
out-of-memory retries in the \ui{LOG}, stay in English in every language, as does the command
line.

\subsection{The Help menu}
\label{sec:settings-help}

\begin{description}[leftmargin=0.6cm,style=nextline]
  \item[\menu{Help}{Models…}]
    The Models window: which models are installed, downloads, import of model files for an
    offline computer, and the model folder (Section~\ref{sec:installing-models}). It also opens by
    itself when a run needs a model that is missing and cannot be downloaded.
  \item[\menu{Help}{Self-test}]
    Runs every installed model on a synthetic image pair with a known translation and reports
    whether each one finds it (Section~\ref{sec:installing-check}). Use it after an installation,
    an upgrade or a change of PyTorch.
  \item[\menu{Help}{About RAFTcorr}]
    The version number, the four models and the displacement convention: reference to deformed,
    $+u$ right, $+v$ down, in pixels, in the coordinates of frame~1. Give this version number when
    you report a problem (Section~\ref{sec:troubleshooting-report}).
\end{description}

\subsection{The File menu}
\label{sec:settings-file}

\begin{description}[leftmargin=0.6cm,style=nextline]
  \item[\menu{File}{Open Image Folder…}]
    Loads a folder of images, as the drop zone of the \ui{IMAGES} section does
    (Section~\ref{sec:images}).
  \item[\menu{File}{Open Project…}, \ui{Save Project}, \ui{Save Project As…}]
    Projects, \file{.raftcorr} files that hold the images, the regions of interest, the settings
    and optionally the results (Section~\ref{sec:projects}).
  \item[\menu{File}{Export Results…}]
    The export dialog, as the \ui{Export Results} button opens it (Section~\ref{sec:export}). It
    does nothing while there are no results.
  \item[\menu{File}{Associate .raftcorr files with RAFTcorr…}]
    Windows only: makes a double-click on a project file open this copy of RAFTcorr
    (Section~\ref{sec:projects-doubleclick}).
  \item[\menu{File}{Quit}]
    Closes RAFTcorr, as closing the main window does. A run in progress is asked to stop first,
    the strain window closes, and the temporary folder of large results is deleted.
\end{description}

\begin{warnbox}
RAFTcorr does not ask before it closes. Regions of interest, results and probes that are not in a
saved project are lost: save a project first (Section~\ref{sec:projects}).
\end{warnbox}

\subsection{Keyboard shortcuts}
\label{sec:settings-keys}

Table~\ref{tab:shortcuts} lists the shortcuts of the menus and the other keys the application
answers to. On macOS, press \kbd{Cmd} where the table says \kbd{Ctrl}.

\begin{table}[htbp]
\centering
\small
\caption{Keyboard shortcuts.}
\label{tab:shortcuts}
\begin{tabularx}{\linewidth}{@{}>{\raggedright\arraybackslash}p{0.3\linewidth}Xl@{}}
  \toprule
  Keys & Action & Section \\
  \midrule
  \kbd{Ctrl+I} & \menu{File}{Open Image Folder…} & \ref{sec:images} \\
  \kbd{Ctrl+O} & \menu{File}{Open Project…} & \ref{sec:projects} \\
  \kbd{Ctrl+S} & \menu{File}{Save Project} & \ref{sec:projects} \\
  \kbd{Ctrl+Shift+S} & \menu{File}{Save Project As…} & \ref{sec:projects} \\
  \kbd{Ctrl+E} & \menu{File}{Export Results…} & \ref{sec:export} \\
  \kbd{Cmd+Q} & \menu{File}{Quit} on macOS (\kbd{Ctrl+Q} on Linux). Windows has no shortcut
    for it; close the window with \kbd{Alt+F4} & \\
  \midrule
  \kbd{Up}, \kbd{Down}, \kbd{Page Up}, \kbd{Page Down}, \kbd{Home}, \kbd{End}
    & In the image list: move to another frame & \ref{sec:images} \\
  \kbd{Esc} & On the canvas: cancel the shape being drawn & \ref{sec:roi} \\
  \kbd{Esc} & In the \ui{Analysis} tab: cancel the placement of a probe, or deselect & \ref{sec:analysis} \\
  \kbd{Delete} & In the \ui{Analysis} tab: delete the selected probe & \ref{sec:analysis} \\
  \kbd{F2} & In the \ui{Analysis} tab: rename the selected probe & \ref{sec:analysis} \\
  \bottomrule
\end{tabularx}
\end{table}
